\begin{examplebox}
\textbf{\textcolor{CExample}{例 1（基础）}}

\textbf{\textcolor{CExample}{题目：}} 椭圆 $\dfrac{x^2}{4}+\dfrac{y^2}{3}=1$，直线 $y=2x+1$，判断直线与椭圆是否相切。

\textbf{\textcolor{CExample}{解题步骤：}}
\begin{description}[style=nextline, leftmargin=2.8em, labelsep=0.5em, itemsep=0.45em]
    \item[\textbf{第一步（提取参数）：}] 从椭圆方程得 $a^2=4$, $b^2=3$；从直线得 $k=2$, $m=1$。
    \par\textbf{理由：} 识别公式 $m^2 = a^2 k^2 + b^2$ 所需参数。

    \item[\textbf{第二步（计算左边）：}] 计算 $m^2 = 1^2 = 1$。
    \par\textbf{理由：} 公式左边是截距平方。

    \item[\textbf{第三步（计算右边）：}] 计算 $a^2 k^2 + b^2 = 4 \times 2^2 + 3 = 4 \times 4 + 3 = 16 + 3 = 19$。
    \par\textbf{理由：} 代入公式右边表达式。

    \item[\textbf{第四步（比较判断）：}] 因为 $1 \neq 19$，所以 $m^2 \neq a^2 k^2 + b^2$。
    \par\textbf{理由：} 充要条件不满足，故直线与椭圆不相切。
\end{description}

\textbf{\textcolor{CExample}{关键结论：}} \textbf{直接应用公式快速判断：若 $m^2 = a^2 k^2 + b^2$ 成立则相切，否则不相切。}

\medskip
\textbf{\textcolor{CExample}{例 2（稍变形）}}

\textbf{\textcolor{CExample}{题目：}} 椭圆 $\dfrac{x^2}{9}+\dfrac{y^2}{5}=1$，直线 $y=kx+2$ 与椭圆相切，求斜率 $k$ 的值。

\textbf{\textcolor{CExample}{解题步骤：}}
\begin{description}[style=nextline, leftmargin=2.8em, labelsep=0.5em, itemsep=0.45em]
    \item[\textbf{第一步（代入公式）：}] 由相切条件 $m^2 = a^2 k^2 + b^2$，代入 $a^2=9$, $b^2=5$, $m=2$。
    \[
    2^2 = 9 k^2 + 5
    \]
    \par\textbf{理由：} 已知相切，直接套用结论建立方程。

    \item[\textbf{第二步（解方程）：}] 计算 $4 = 9k^2 + 5$，移项得 $9k^2 = -1$。
    \par\textbf{理由：} 整理方程求解 $k^2$。

    \item[\textbf{第三步（分析结果）：}] $9k^2 = -1$ 无实数解，因为平方数非负。
    \par\textbf{理由：} 实数范围内 $k^2 \ge 0$，右边为负矛盾。

    \item[\textbf{第四步（结论）：}] 不存在实数 $k$ 使直线与椭圆相切。
    \par\textbf{理由：} 方程无解意味着给定截距 $m=2$ 下，无法找到斜率 $k$ 使直线相切。
\end{description}

\textbf{\textcolor{CExample}{关键结论：}} \textbf{利用公式列方程时，注意实数解的存在性；本例显示截距 $m=2$ 过大，无法与椭圆相切。}
\end{examplebox}
